\section{Off-policy and on-policy methods}
The convergence of RL with the substantial representational capabilities of deep neural networks has given rise to \textbf{Deep Reinforcement Learning (DRL)}, which currently represents the leading edge of V2G control research. The development of DRL algorithms has yielded a comprehensive toolkit, primarily divided into two main families: off-policy and on-policy methods, each exhibiting distinct operational characteristics.
The following figure illustrates the evolutionary relationships between the algorithms that  will be discussed, highlighting how newer ideas were built to overcome the limitations of their predecessors.
\\
\noindent
\textit{It can be considered as a \textbf{rule of thumb} in the case of actor-critic algorithms that: many off-policy algorithms use the critic as a Q-function because it allows updating the actor using experiences from the replay buffer, without strictly following the current policy (e.g., DDPG, SAC, TQC); while on-policy algorithms often use the critic as a V-function because directly evaluating the state with the current policy is more stable (e.g., A2C, PPO).}

\begin{figure}[H]
\centering
\begin{forest}
for tree={
  grow=east,
  parent anchor=east,
  child anchor=west,
  edge path={
    \noexpand\path [draw, thick, \forestoption{edge}]
      (!u.parent anchor) -- +(-0.5pt,0) |- (.child anchor);
  },
  draw,
  rounded corners,
  align=center,
  l sep=8pt,
  s sep=6pt,
  minimum height=0.5cm,
  inner sep=2pt,
  anchor=west,
  font=\sffamily\small,
}
[RL Methods, fill=blue!20, edge=blue!60
    [On-Policy, fill=green!20, edge=green!60
        [Policy Gradient, fill=green!10, edge=green!40
            [A2C, fill=yellow!15, edge=green!50]
            [TRPO\\(Theoretical Stability), fill=yellow!15, edge=green!50]
            [PPO\\(Practical Simplification), fill=yellow!15, edge=green!50]
        ]
        [Gradient-Free, fill=teal!15, edge=teal!50
            [ARS\\(Random Search), fill=yellow!15, edge=teal!50]
        ]
    ]
    [Off-Policy, fill=red!20, edge=red!60
        [DDPG\\(Fundamental), fill=red!10, edge=red!50
            [DDPG+PER\\(Improve Sampling), fill=orange!15, edge=red!50]
            [TD3\\(Resolve Overestimation), fill=orange!15, edge=red!50]
            [SAC\\(Add Entropy), fill=orange!15, edge=red!50]
            [TQC\\(Distributional Approach), fill=orange!15, edge=red!50]
        ]
    ]
]
\end{forest}
\caption{Hierarchical classification of Reinforcement Learning methods.}
\label{fig:rl_methods}
\end{figure}

\subsection{Off-policy Reinforcement Learning}

Off-policy algorithms distinguish themselves through their ability to learn optimal policies from data generated by different, often more exploratory, behavioral policies. \textbf{This enables them to reuse historical experiences stored in large \textit{replay buffers}}, breaking temporal correlations between experiences and achieving high sample efficiency. A central challenge in reinforcement learning is balancing the exploration of new actions with the exploitation of known good actions. Off-policy methods are particularly well-suited to this task, as they can leverage past experiences to guide learning without being constrained by the current policy.

\subsubsection{Soft Actor-Critic (SAC)}

Soft Actor-Critic (SAC) \cite{haarnoja2018soft,haarnoja2019soft} is a state-of-the-art off-policy, model-free, actor-critic algorithm designed for continuous control tasks. It is recognized for its \textbf{superior sample efficiency and stability}. SAC's key innovation lies in its \textbf{maximum entropy framework}, where the agent's objective is to maximize both expected reward and policy entropy. This entropy bonus encourages agents to explore broadly, improving noise robustness and reducing the risk of converging to poor local optima. 
\noindent
By explicitly promoting randomness in the policy, SAC addresses the exploration-exploitation trade-off in a principled manner. \cite{gu2025optimization} extend this idea to safe reinforcement learning, considering constraints such as battery degradation to maintain reliable and safe operations in real-world Vehicle-to-Grid (V2G) applications.

\subsubsection{Deep Deterministic Policy Gradient (DDPG)}

Deep Deterministic Policy Gradient (DDPG) \cite{lillicrap2015continuous} extends Deep Q-Networks (DQN) to high-dimensional, continuous action spaces, providing a foundational actor-critic approach for control problems, including V2G. The actor deterministically maps states to actions, while the critic estimates Q-values. However, DDPG suffers from practical challenges such as training instability and \textbf{overestimation bias}, where critics systematically overestimate Q-values, leading to suboptimal policy learning \cite{orfanoudakis2024ev2gym, alfaverh2022optima}. Despite these limitations, DDPG has been successfully applied in V2G contexts, including demand response management \cite{wang2022electrical} and large-scale coordination through transfer learning \cite{Zhang2023}.


\subsubsection{Enhancement: Prioritized Experience Replay (PER)}

This represents not a standalone algorithm but a crucial orthogonal modification for off-policy methods. Standard replay buffers sample past transitions uniformly. PER samples transitions based on their "importance," typically proportional to TD error magnitude \cite{schaul2015prioritized}. This focuses learning processes on surprising or informative experiences, significantly accelerating convergence and improving final performance.


\subsubsection{Truncated Quantile Critics (TQC)}

Truncated Quantile Critics (TQC) \cite{kuznetsov2020controlling} addresses overestimation bias through a distributional RL approach. Instead of learning only expected returns, TQC estimates \textbf{full return distributions using quantile\footnote{A quantile is a value that divides a data distribution into equal parts.} regression}. 

\noindent
Quantile regression extends classical linear regression by estimating conditional quantiles rather than only the mean. Formally, given data $\{(x_i, y_i)\}_{i=1}^n$ with predictors $x_i \in \mathbb{R}^p$ and response $y_i \in \mathbb{R}$, the $\tau$-th conditional quantile ($0 < \tau < 1$) of $Y$ given $X=x$ is:
\[
Q_Y(\tau \mid X=x) = x^\top \beta_\tau,
\]
where $\beta_\tau$ are the regression coefficients associated with quantile $\tau$. The estimator is obtained by solving:
\[
\hat{\beta}_\tau = \underset{\beta \in \mathbb{R}^p}{\arg\min} 
\sum_{i=1}^n \rho_\tau \big( y_i - x_i^\top \beta \big),
\]
with the asymmetric loss function $\rho_\tau(u) = u \cdot (\tau - \mathbf{1}_{\{u < 0\}})$. For example, $\tau=0.5$ corresponds to the conditional median.

\noindent
By learning multiple quantile estimates and truncating the most optimistic ones before averaging, TQC \textbf{reduces overestimation bias} effectively, leading to improved performance in continuous control tasks while maintaining robustness.


\subsubsection{Twin Delayed DDPG (TD3)}

TD3 \cite{fujimoto2018addressing} was developed to mitigate DDPG's instabilities and overestimation bias. It introduces three main innovations: 
(1) \textbf{clipped double Q-learning}, using two critic networks and taking the minimum estimate; 
(2) \textbf{delayed policy updates}, updating the actor less frequently than the critics; 
(3) \textbf{target policy smoothing}, adding noise to target actions for regularization. 

These enhancements make TD3 a robust and reliable baseline for complex V2G scheduling tasks. Studies have applied TD3 to optimize EV charging and discharging schedules while considering grid stability \cite{Dou, Ding}, demonstrating its ability to balance economic incentives with physical constraints.

\subsection{On-Policy Methods: Stability through Cautious Updates}

On-policy methods learn exclusively from data generated by current policies being optimized. Once data is used for updates, it is discarded. While this makes them inherently less sample-efficient than off-policy counterparts, their updates often demonstrate greater stability and reduced divergence risk.

\subsubsection{Advantage Actor-Critic (A2C/A3C)}

A2C represents a foundational synchronous, on-policy Actor-Critic algorithm. Its powerful extension, \textbf{Asynchronous Advantage Actor-Critic (A3C)}, was a landmark contribution demonstrating parallelism benefits \cite{mnih2016asynchronous}. A3C utilizes multiple parallel workers, each with individual model and environment copies. These workers interact with their environments independently, and collected gradients update a global model asynchronously. This decorrelates data streams and provides powerful stabilizing effects on learning processes.
\\
\noindent
While often considered less sample-efficient than their off-policy counterparts, on-policy methods like A2C and its variants are still actively explored in the V2G domain. For example, Chifu et al. \cite{chifu2024deep} propose a Deep Q-Learning based approach for smart scheduling of EVs for demand response, which shares some of the on-policy characteristics. Their work raises the question:
\\
\noindent
Can the stability and reliability of on-policy training outweigh the sample efficiency of off-policy methods in a problem where the environment is highly stochastic? This thesis provides a direct comparison by benchmarking several state-of-the-art off-policy agents against on-policy baselines, allowing for an empirical test of this long-standing hypothesis in the DRL community.

\subsubsection{Trust Region Policy Optimization (TRPO)}

TRPO was the first algorithm to formalize policy update size constraints for guaranteed monotonic policy improvement \cite{schulman2020trust}. It maximizes a "surrogate" objective function subject to policy change constraints, measured by Kullback-Leibler (KL) divergence. This creates a "trust region" within which new policies are guaranteed to improve upon previous ones, preventing catastrophic updates that can permanently derail learning. However, implementation complexity arises from second-order optimization requirements.

\subsubsection{Proximal Policy Optimization (PPO)}

PPO achieves TRPO's stability benefits and reliable performance using only first-order optimization, making it far simpler to implement and more broadly applicable \cite{schulman2017proximal}. It employs a novel \textbf{clipping} mechanism in its objective function to discourage large policy updates that would move new policies too far from previous ones, effectively creating "soft" trust regions. Due to its excellent balance of performance, stability, and implementation simplicity, PPO has become a default choice for many on-policy applications.

\subsection{Gradient-Free Methods: An Alternative Path}

\subsubsection{Augmented Random Search (ARS)}

As a counterpoint to dominant gradient-based methods, ARS represents a gradient-free approach that optimizes policies directly in parameter space \cite{mania2018simple}. It operates by exploring random policy parameter perturbations and updating central policy parameter vectors based on observed performance of these perturbations. While often less sample-efficient for complex, high-dimensional problems, its extreme simplicity, scalability, and robustness to noisy rewards can make it competitive in certain domains.
